\documentclass[12pt]{article}
\usepackage[utf8]{inputenc}
\usepackage[T1]{fontenc}
\usepackage{hyperref}
\title{Xamtanga Language Resources for K-12 Curriculum Development}
\author{Compiled by Claude Sonnet 5.5 (claude-sonnet-5-5) for Andr\'as Kornai}
\date{}
\begin{document}
\maketitle

\section{Introduction}
Xamtanga (ISO 639-3 xan, Glottocode xamt1239) is a Central Cushitic (Agaw) language of the northern Ethiopian highlands. It is spoken by the Xamir people, mainly in the Wag Hemra Zone of the Amhara Region. About 210,000 first-language speakers are reported, a figure from 2010. Bilingualism in Amharic or Tigrinya is very common. One source describes Xamtanga as a working and educational language in Wag Hemra, but these searches found no data on which schools teach it. No primers or textbooks were found.

\section{Primary Educational Resources}
None were found. A Seed Company project profile notes that literacy is low among the Wag Himra people: \url{https://seedcompany.com/wp-content/uploads/2024/07/Xamtanga-3-Project-Profile.pdf}

\section{Linguistic and Reference Materials}
\begin{itemize}
\item Xamtanga language overview (Wikipedia): \url{https://en.wikipedia.org/wiki/Xamtanga_language}
\item Chlo\'e Darmon (2015), \emph{A morphosyntactic description of Xamtanga}, Lyon 2 doctoral dissertation. Glottolog record: \url{https://glottolog.org/resource/reference/id/548296}
\item Darmon's chapter on Xamtanga in \emph{The Oxford Handbook of Ethiopian Languages}: \url{https://academic.oup.com/edited-volume/46265/chapter/405502216}
\item Darmon, ``The morpheme -(\"a)\~n\"a in Xamtanga'': \url{https://benjamins.com/catalog/tsl.117.15dar}
\item David Appleyard, \emph{A comparative dictionary of the Agaw languages} (K\"oln: R\"udiger K\"oppe, 2006). It covers Agaw generally, and no URL for it was seen.
\item David Appleyard, ``A Grammatical Sketch of Khamtanga'', \emph{Bulletin of SOAS} (1987). No URL for it was seen.
\item ASJP wordlist: \url{https://asjp.clld.org/languages/XAMTANGA}
\end{itemize}

\section{Religious and Community Resources}
The Seed Company profile says a New Testament exists in Xamtanga and that Old Testament translation is underway. No direct URL for the scripture itself was seen. Background on the Xamir people: \url{https://en.wikipedia.org/wiki/Xamir_people}

\section{Digital and Audio Resources}
No audio, video or software resources were found.

\section{Curriculum Development Context}
Ethiopia introduced mother-tongue primary education in 1994. Searches found no policy or programme specific to Xamtanga. The Wag Hemra Zone background is at \url{https://en.wikipedia.org/wiki/Wag_Hemra_Zone}. Any curriculum would have to start from the linguistic descriptions above and the Bible translation work, since no school materials were found.

\end{document}